\section{From curve defects to matrix codimension}
\label{sec:dimension}

The preceding sections show that a critical matrix of the coefficient
polynomial produces a reduced plane curve with large normalization defect.
We now count these matrices by separating the choice of the curve from the
choice of a matrix that produces it.

\begin{samepage}
\begin{proposition}[Gradient codimension]\label{prop:gradient-codimension}
Let $K=\overline{\mathbb F}_2$, let $k\geq2$, and set
\[
 q=\left\lfloor\frac{k-1}{3}\right\rfloor,
 \qquad T=\binom{q+2}{2}.
\]
For pairwise distinct $\theta_1,\ldots,\theta_k\in K$, put
$D=\operatorname{diag}(\theta_1,\ldots,\theta_k)$ and
\[
 G_D(X)=[s^qt^q]\per_k(X+sI+tD),
 \qquad X\in\Mat_k(K).
\]
Then
\begin{equation}
 \dim\{X\in\Mat_k(K):\nabla G_D(X)=0\}
 \leq k^2+k-T.
 \label{eq:gradient-dimension}
\end{equation}
\end{proposition}
\end{samepage}

For the proof, fix these parameters and write
\[
 \Gamma_D=\{X\in\Mat_k(K):\nabla G_D(X)=0\}.
\]
Consider the morphism
\begin{equation}
 \label{eq:determinant-curve-map}
 f_D:\Mat_k(K)\longrightarrow
 \PP H^0(\PP^2_K,\cO(k)),\qquad
 X\longmapsto[F_X],\qquad
 F_X=\det(sI+tD+uX).
\end{equation}
The coefficient of $s^k$ is one, so a projective fiber fixes the
polynomial $F_X$ itself and hence the embedded curve $C_X$.
Every curve in the image is reduced by
Lemma~\ref{lem:adjugate-projection}.  For $X\in\Gamma_D$,
Propositions~\ref{prop:adjugate-length} and~\ref{prop:conductor-length}
give
\begin{equation}
 \label{eq:dimension-gradient-defect}
 T\leq\length Z\leq\delta(C_X).
\end{equation}
We will bound the dimension of $f_D(\Gamma_D)$ using this defect, then
bound the fibers of $f_D$ by associating a finite subscheme of the fixed
curve to each matrix in the fiber.

\subsection{Curves with large normalization defect}

Recall that, for a reduced curve $C$ with normalization
$\nu:\widetilde C\to C$,
\[
 \delta(C)=\length(\nu_*\cO_{\widetilde C}/\cO_C).
\]
The normalization may be disconnected.  We will keep track of its
components explicitly when applying the dimension theorem for integral
curves.

Over any algebraically closed field, let $b\geq1$ and
$0\leq g\leq(b-1)(b-2)/2$ be integers.  The locus of integral plane
curves of degree $b$ and normalization genus $g$ is locally closed and,
when nonempty, has dimension $3b+g-1$
\cite[Lemma~2.6, Proposition~2.7, and Remark~2.8]{CHT2023}.
This is the cited theorem with surface $\PP^2$ and line bundle $\cO(b)$.
The locus includes every integral curve of the stated genus: the surface
has no singular points to exclude, and nodality is not a hypothesis.

\begin{lemma}
Let $K$ be an algebraically closed field\label{lem:severi-dimension}, and let $k\geq1$ and $T\geq0$
be integers.  In the degree-$k$ plane linear system over $K$, the locus
of reduced curves satisfying $\delta(C)\geq T$ is constructible.
If nonempty, it has dimension at most
\[
 \frac{k(k+3)}2-T.
\]
\end{lemma}

\begin{proof}
Write $C=C_1\cup\cdots\cup C_a$, where the integral components have
degrees $b_i$ and normalization genera $g_i$.  The Euler characteristics
of the normalization and the curve are $a-\sum_i g_i$ and $1-p_a(C)$.
The normalization exact sequence therefore gives
\[
 \delta(C)=p_a(C)-\sum_{i=1}^a g_i+a-1,
 \qquad p_a(C)=\frac{(k-1)(k-2)}2.
\]
In particular, the term $a-1$ accounts for the disconnected normalization,
and the identity includes the contribution of intersections between
different components.

For fixed degrees and genera, multiplying the equations of the
components defines a morphism from the product of their integral curve
loci to the degree-$k$ linear system.  Restricting to distinct factors
gives precisely the reduced curves of this type.  The image is
constructible, and its dimension is at most
\begin{align*}
 \sum_{i=1}^a(3b_i+g_i-1)
 &=3k+\sum_{i=1}^a g_i-a\\
 &=\frac{k(k+3)}2-\delta(C).
\end{align*}
There are only finitely many degree and genus types:
$a\leq k$, $\sum b_i=k$, and
$0\leq g_i\leq(b_i-1)(b_i-2)/2$.
The normalization identity determines $\delta(C)$ from each type.
The required locus is therefore a finite union of constructible images
with the stated dimension bound.
\end{proof}

\subsection{A row subscheme determines a matrix}

Fix a curve $C$ in the image of $f_D$, and let
$\mathcal F_C=f_D^{-1}([C])$.  Our remaining task is to bound this
matrix fiber.  For a critical matrix, its full-adjugate scheme $Z$
supplied the defect bound \eqref{eq:dimension-gradient-defect} on its
curve.  We now use a row obtained from the adjugate to distinguish all
matrices in the fixed fiber $\mathcal F_C$.

Put $N=k(k-1)/2$.  For $X\in\mathcal F_C$, define
\[
 R_X=(1,\ldots,1)\adj(sI+tD+uX),
\]
and define an ideal and its subscheme on the fixed curve by
\begin{equation}
 \label{eq:row-adjugate-ideal}
 J'_X=\operatorname{im}\bigl(
 \cO_C(1-k)^k\xrightarrow{\ R_X\ }\cO_C\bigr),
 \qquad Z'_X=V_C(J'_X).
\end{equation}
Thus $J'_X$ is locally generated by the row entries after trivializing
their common twist.  We use the prime to distinguish this row subscheme
from the full-adjugate scheme $Z$.

After transposition, the equations $R_X=0$ say that an adjugate matrix
annihilates a fixed vector.  Adams, Harnad and Hurtubise use equations
of this form for spectral divisors in a generic complex setting
\cite[Section~1b, equation~(1.29) in arXiv:hep-th/9210089v1]{AdamsHarnadHurtubise1993}.
Here we prove the uniform finite length of $Z'_X$ directly.  The common
boundary frame then makes this subscheme determine the matrix uniquely
on a fixed embedded curve.

\begin{lemma}[The row subscheme]
For each $X\in\mathcal F_C$\label{lem:row-adjugate-scheme}, the scheme $Z'_X$ is finite of length $N$.
If $X,X'\in\mathcal F_C$ satisfy $Z'_X=Z'_{X'}$ as subschemes of $C$,
then $X=X'$.
\end{lemma}

\begin{proof}
The projection $\pi:C\to\PP^1_{t:u}$ and the bundle
\[
 \mathcal E=\pi_*\cO_C(k-1)
 =\bigoplus_{j=0}^{k-1}\cO(k-1-j)
\]
are fixed throughout $\mathcal F_C$.
The entries of $R_X$ give a map
$\rho_X:\cO^k\to\mathcal E$ on $\PP^1$.
At $b_0=[1:0]$, trivialize by $t$ and put $y=s/t$.  For every $X$,
the map on this fiber sends the $i$-th standard basis vector to
\begin{equation}
 \label{eq:row-boundary-frame}
 h_i(y)=\prod_{j\ne i}(y+\theta_j)
 \quad\text{in}\quad
 K[y]\big/\prod_j(y+\theta_j).
\end{equation}
Evaluation at the distinct roots $-\theta_i$ shows that these $k$
classes form a basis.  In particular, all the maps $\rho_X$ restrict
to the same isomorphism at $b_0$.  Each $\rho_X$ is therefore
generically full rank and injective as a sheaf map.  Its cokernel is
torsion of length
\begin{equation}
 \label{eq:row-adjugate-length}
 \length(\coker\rho_X)
 =\deg\mathcal E
 =\sum_{j=0}^{k-1}(k-1-j)=N.
\end{equation}

On $C$, the adjugate identity gives
\[
 sR_X=-R_X(tD+uX).
\]
On either affine base chart, the span of the row entries over the base
ring is consequently stable under the affine $s$-coordinate.  This
coordinate generates the finite curve algebra over the base, so the
span is already the ideal generated over the curve algebra, with its
twist.  Hence
\[
 \operatorname{im}\rho_X=\pi_*(J'_X(k-1)).
\]
Finite pushforward is exact, so
\[
 \coker\rho_X=\pi_*(\cO_{Z'_X}(k-1)).
\]
This pushforward has finite support, and $\pi$ is finite, so
$\cO_{Z'_X}$ has finite support as well.  Twisting preserves length,
as does finite pushforward over the algebraically closed ground field.
Equation \eqref{eq:row-adjugate-length} thus proves that $Z'_X$ is
finite of length $N$.

Suppose now that $Z'_X=Z'_{X'}$.  Their ideal sheaves on the fixed
curve are equal, so $\rho_X$ and $\rho_{X'}$ have the same image inside
the same bundle $\mathcal E$.  Since both maps are injective, there is
an automorphism $U$ of $\cO^k$ with
\[
 \rho_{X'}=\rho_XU.
\]
Every such automorphism is an invertible constant matrix.  Restricting
this equality to $b_0$, where both maps are the same isomorphism
\eqref{eq:row-boundary-frame}, gives $U=I$ and hence
$\rho_{X'}=\rho_X$.

Multiplication by $s$ maps $\mathcal E$ to $\mathcal E(1)$, and the
row identity says
\[
 s\rho_X=-\rho_X(1)(tD+uX).
\]
The analogous equation for $X'$ now gives
$\rho_X(1)u(X-X')=0$.  The map $\rho_X(1)$ is injective, so
$u(X-X')=0$ as a matrix of sections of $\cO(1)$.  Since $u$ is a
nonzero section and $X-X'$ is constant, this forces $X=X'$.
The same argument works after extending the ground field, so the
uniqueness statement also holds for geometric points.
\end{proof}

The lemma reduces the fiber count to the possible length-$N$ subschemes
of $C$.  We next bound the dimension of their Hilbert scheme in a form
valid in characteristic two.

\subsection{Finite subschemes of a plane curve}

For a punctual quotient of $\cO$, Baranovsky's construction of punctual
Quot schemes over $\C$ uses commuting nilpotent operators and a cyclic
vector after choosing a basis \cite[Section~2, Lemma~2.2]{Baranovsky1998}.
For the classical punctual-plane Hilbert dimension bound in arbitrary
characteristic, see \cite[Corollary~4.1]{Premet2003}.
The centralizer count below supplies the dimension bound needed here in
every characteristic.

\begin{lemma}\label{lem:hilbert-curve-dimension}
Let $C\subset\PP^2_K$ be a reduced curve over an algebraically closed
field $K$, and let $N\geq0$.  Then
\[
 \dim\operatorname{Hilb}^N(C)\leq N.
\]
\end{lemma}

\begin{proof}
Fix a point of the affine plane and an integer $\ell\geq1$, and
translate the point to the origin.  Inside
$\operatorname{Hilb}^{\ell}(\mathbb A^2_K)$, let $P_\ell$ be the closed
locus of quotients on which $(x,y)^\ell$ acts by zero.
Its geometric points are precisely the punctual subschemes supported at
the origin: in a length-$\ell$ local algebra, the $\ell$-th power of
the maximal ideal is zero.

The frame bundle of the tautological quotient algebra over $P_\ell$
has dimension $\dim P_\ell+\ell^2$.  In a frame, multiplication by
the two coordinates gives commuting nilpotent matrices $A,B$ of size
$\ell$, and the unit gives a cyclic vector $v\in K^\ell$, meaning
$K[A,B]v=K^\ell$.  These data recover the ideal as
\[
 \{f\in K[x,y]:f(A,B)v=0\},
\]
and recover its frame.  Thus the resulting morphism from the frame
bundle to the finite-type scheme of such triples is injective on geometric points.
Its geometric fibers have dimension zero, which suffices for a
dimension upper bound.

There are finitely many nilpotent Jordan types for $A$.  If its linear
centralizer has dimension $j$, then its group centralizer is the open
set of invertible elements in that vector space.  The conjugacy orbit
of $A$ therefore has dimension $\ell^2-j$.  Within its centralizer,
the nilpotent choices of $B$ have dimension at most $j-1$: they satisfy
$\det B=0$, whereas the determinant is one at the identity matrix in
the same centralizer.  This argument works in every characteristic.
Consequently the commuting nilpotent pairs have dimension at most
$\ell^2-1$.  The choice of the cyclic vector adds at most $\ell$
parameters, so
\[
 \dim P_\ell+\ell^2\leq\ell^2+\ell-1,
 \qquad \dim P_\ell\leq\ell-1.
\]

Fix $p\in C$ and choose an affine plane chart around it, translated so
that $p$ is the origin.  Inside the Hilbert scheme of this affine part
of $C$, impose the same condition that the $\ell$-th power of the
maximal ideal of $p$ acts by zero.  The resulting locus is a closed
subscheme of $P_\ell$, and its geometric points are exactly the
length-$\ell$ subschemes of $C$ supported at $p$.  Its dimension is
therefore at most $\ell-1$, including when $p$ is singular.
Fix a positive length type
$\ell_1+\cdots+\ell_a=N$.  Consider ordered distinct points
$(p_1,\ldots,p_a)$ of $C$ together with subschemes of lengths $\ell_i$
supported at $p_i$.  This incidence is algebraic: containment in the
$\ell_i$-th infinitesimal neighborhood of the moving point can be
imposed using the $\ell_i$-th power of the diagonal ideal on $C\times C$.
The support parameters have dimension at most $a$, and the fibers have
dimension at most $\sum_i(\ell_i-1)=N-a$.  The incidence therefore
has dimension at most $N$.  Taking disjoint unions of its subschemes
maps it to $\operatorname{Hilb}^N(C)$.  Finitely many length types
cover all geometric points of this Hilbert scheme, proving the claim.
The case $N=0$ is immediate.
\end{proof}

\subsection{The matrix fibers and the gradient locus}

We now use the universal family of length-$N$ subschemes to compare the
matrix fiber with $\operatorname{Hilb}^N(C)$.

\begin{proposition}[Dimension of a matrix fiber]\label{prop:matrix-fiber}
For the morphism $f_D$ in \eqref{eq:determinant-curve-map}, every fiber
has dimension at most $N=k(k-1)/2$.
\end{proposition}

\begin{proof}
Fix a curve $C$ in the image and its fiber $\mathcal F_C$.  Form
\[
 \mathcal I_C=
 \{(X,Z')\in\mathcal F_C\times\operatorname{Hilb}^N(C):
 R_X|_{Z'}=0\}.
\]
This is a closed incidence.  Indeed, restrict the row entries to the
universal finite flat subscheme over the Hilbert scheme and push forward
their common twist.  This produces sections of a vector bundle on the
product, whose zero locus is $\mathcal I_C$.

Every matrix $X$ occurs in the incidence with its subscheme $Z'_X$.
Conversely, row vanishing on $Z'$ means $J'_X\subset I_{Z'}$, so
$Z'\subset Z'_X$.  The scheme $Z'$ has length $N$ by its membership
in the Hilbert scheme, and $Z'_X$ has length $N$ by
Lemma~\ref{lem:row-adjugate-scheme}.  Therefore they are equal.
The same lemma shows that each fiber of $\mathcal I_C$ over
$\operatorname{Hilb}^N(C)$ has at most one geometric point.  Such a
finite-type fiber has dimension zero when nonempty, even if it is
nonreduced.  Projection to $\mathcal F_C$ is surjective on geometric
points, so
\[
 \dim\mathcal F_C
 \leq\dim\mathcal I_C
 \leq\dim\operatorname{Hilb}^N(C)
 \leq N,
\]
where the last inequality is Lemma~\ref{lem:hilbert-curve-dimension}.
\end{proof}

\begin{proof}[Proof of Proposition~\ref{prop:gradient-codimension}]
By \eqref{eq:dimension-gradient-defect}, every curve in
$f_D(\Gamma_D)$ has normalization defect at least $T$.
Lemma~\ref{lem:severi-dimension} bounds the dimension of this
constructible image by $k(k+3)/2-T$.  Each fiber of the restriction
of $f_D$ to $\Gamma_D$ is contained in a full matrix fiber and has
dimension at most $N$ by Proposition~\ref{prop:matrix-fiber}.
Consequently
\[
 \dim\Gamma_D
 \leq\frac{k(k+3)}2-T+\frac{k(k-1)}2
 =k^2+k-T,
\]
as asserted.
\end{proof}
